\documentclass[../DoAn.tex]{subfiles}
\begin{document}


Chương 1 đặt nền móng cho toàn bộ nghiên cứu của đồ án bằng cách giới thiệu bối cảnh bài toán, đặt vấn đề khoa học và thực tiễn, từ đó xác lập mục tiêu, phạm vi nghiên cứu cùng định hướng giải pháp kỹ thuật cụ thể. Nội dung chương sẽ lần lượt trình bày về: đặt vấn đề về sự bùng nổ của mô hình thương mại điện tử đa nhà bán hàng cùng các thách thức lớn về tải và bảo mật trong phần \ref{section:1.1}; mục tiêu và phạm vi cụ thể của đề tài trong phần \ref{section:1.2}; định hướng giải pháp kỹ thuật lựa chọn bao gồm MERN Stack, Redis và Socket.IO trong phần \ref{section:1.3}; và sau cùng là bố cục chi tiết của quyển báo cáo đồ án tốt nghiệp trong phần \ref{section:1.4}.

\section{Đặt vấn đề}
\label{section:1.1}
Trong kỷ nguyên số hóa mạnh mẽ và sự chuyển dịch kinh tế toàn cầu, thương mại điện tử (E-commerce) đã trở thành một phần tất yếu trong đời sống xã hội và là động lực phát triển kinh tế quan trọng. Đặc biệt, mô hình sàn giao dịch thương mại điện tử đa nhà bán hàng (Multi-vendor Marketplace) – nơi nhiều thương nhân cùng đăng ký mở gian hàng độc lập và tương tác với đông đảo người mua – đang chiếm ưu thế vượt trội nhờ khả năng tối ưu hóa chi phí vận hành và đa dạng hóa nguồn hàng. Tuy nhiên, sự phát triển bùng nổ này cũng đi kèm với những thách thức công nghệ vô cùng phức tạp mà các nền tảng thương mại điện tử thế hệ cũ chưa thể giải quyết triệt để.

Thách thức lớn nhất đến từ sự tranh chấp tài nguyên trong môi trường đồng thời tải cao (High Concurrency). Vào các khung giờ cao điểm hoặc chiến dịch Flash Sale lớn, hàng chục nghìn lượt truy cập đồng thời đổ dồn về một số mặt hàng hot, dẫn đến hiện tượng tranh chấp tài nguyên ở tầng dữ liệu. Các kiến trúc thông thường sử dụng cơ chế đọc-sửa-ghi (Read-Modify-Write) rất dễ dẫn đến lỗi bán vượt quá số lượng thực tế trong kho (overselling), gây thiệt hại đến uy tín của nhà bán hàng và làm gián đoạn trải nghiệm người tiêu dùng. Bên cạnh đó, an toàn thông tin và bảo mật phiên đăng nhập cũng là vấn đề nan giải. Trong bối cảnh sàn đa nhà bán hàng, việc một người dùng có thể đồng thời tham gia với vai trò người mua và người bán đặt ra yêu cầu chặt chẽ về xác thực phiên, kiểm soát truy cập và chống lạm dụng tài nguyên đăng nhập như Brute-force hoặc spam OTP. Ngoài ra, việc thiếu đi các mô hình cá nhân hóa trải nghiệm khách hàng khiến nền tảng giảm tỷ lệ chuyển đổi, trong khi các thuật toán khuyến nghị sản phẩm truyền thống thường chịu hạn chế về tính cập nhật và gặp lỗi khởi đầu lạnh (Cold Start). Xuất phát từ thực tế đó, đề tài được thực hiện nhằm nghiên cứu và xây dựng một hệ thống sàn thương mại điện tử đa nhà bán hàng có khả năng vận hành ổn định dưới tải cao, bảo đảm an toàn cho phiên đăng nhập và nâng cao trải nghiệm mua sắm cá nhân hóa.

\section{Mục tiêu và phạm vi đề tài}
\label{section:1.2}
Hiện nay, đa số các nghiên cứu và sản phẩm thương mại điện tử hiện có trên thị trường thường tập trung vào việc làm phong phú thêm giao diện người dùng mà bỏ qua việc tối ưu hóa tầng xử lý giao dịch đồng thời và an toàn dữ liệu ở tầng backend. Các sàn giao dịch hiện tại thường gánh chịu độ trễ lớn khi thực hiện tính toán các khuyến mãi phức tạp, hoặc phải sử dụng các giải pháp phần cứng đắt đỏ để che giấu các lỗ hổng kiến trúc phần mềm. 

Dựa trên các phân tích và đánh giá trên, đồ án hướng tới mục tiêu xây dựng một nền tảng sàn thương mại điện tử đa nhà bán hàng toàn diện, an toàn, có khả năng chịu tải tốt và hỗ trợ cá nhân hóa trải nghiệm người dùng. Về mặt phạm vi, đề tài tập trung giải quyết các nhiệm vụ cốt lõi sau:
(i) Thiết kế và phát triển hệ thống backend chịu tải cao sử dụng cơ chế giữ kho nguyên tử ở tầng cơ sở dữ liệu kết hợp khóa chống trùng đơn (idempotency key) nhằm kiểm soát chặt chẽ nguy cơ oversell;
(ii) Củng cố bảo mật hệ thống thông qua mô hình tài khoản hợp nhất (Unified User/Vendor), cơ chế giới hạn tần suất truy cập bằng Redis và quản lý phiên xác thực phía server kết hợp xoay vòng refresh token;
(iii) Hiện thực hóa công cụ tính giá và khuyến mãi nhiều lớp tại backend, bao gồm giá sản phẩm, giảm giá theo shop, giảm giá nền tảng và giảm giá phí vận chuyển;
(iv) Xây dựng động cơ khuyến nghị sản phẩm lai (Hybrid Recommendation Engine) dựa trên điểm tương tác người dùng (Interaction Score) có áp dụng hệ số suy hao theo thời gian; và
(v) Thiết kế giao diện React đồng bộ bao gồm trang mua sắm cho người dùng cuối và bảng điều khiển phân tích số liệu trực quan cho đối tác bán hàng.

\section{Định hướng giải pháp}
\label{section:1.3}
Để đạt được các mục tiêu đặt ra, đồ án lựa chọn bộ giải pháp công nghệ hiện đại và đồng bộ sau:
\begin{itemize}
\item \textbf{Kiến trúc MERN Stack (MongoDB, Express, React, Node.js):} Sử dụng Node.js và Express ở backend cung cấp cơ chế non-blocking I/O hiệu năng cao để xử lý hàng nghìn kết nối đồng thời. Cơ sở dữ liệu MongoDB hỗ trợ mô hình tài liệu linh hoạt (document-based), rất phù hợp với cấu trúc phân cấp đa dạng của sản phẩm và đơn hàng đa nhà bán hàng. React mang lại giao diện người dùng mượt mà kiểu Single Page Application (SPA).
\item \textbf{Bộ nhớ đệm tốc độ cao Redis:} Đóng vai trò là lớp giới hạn tần suất truy cập (Rate-limiter) và kho lưu trữ phiên xác thực phía server, phục vụ cho cơ chế xoay vòng refresh token và thu hồi phiên đăng nhập khi cần thiết, đồng thời giúp giảm tải cho cơ sở dữ liệu chính MongoDB.
\item \textbf{Kênh truyền thời gian thực Socket.IO:} Thiết lập kênh truyền dữ liệu hai chiều thời gian thực (duy trì bằng JWT auth), phục vụ cho hệ thống chat trực tuyến tức thời giữa người mua và người bán.
\item \textbf{Tích hợp thanh toán trực tuyến qua cổng Stripe và VNPay:} Bảo đảm luồng giao dịch tài chính diễn ra an toàn và minh bạch, trong đó Stripe sử dụng webhook để xác nhận thanh toán còn VNPay sử dụng cơ chế callback/redirect có ký số để đối soát kết quả giao dịch.
\end{itemize}

\section{Bố cục đồ án}
\label{section:1.4}
Phần còn lại của báo cáo đồ án tốt nghiệp này được tổ chức như sau.

Chương 2 trình bày về phân tích và thiết kế hệ thống bao gồm khảo sát hiện trạng thực tế, mô tả tổng quan chức năng qua các biểu đồ ca sử dụng (Use Case) tổng quát và phân rã các nhóm nghiệp vụ chính, mô phỏng quy trình nghiệp vụ đặt hàng giữ kho nguyên tử kết hợp thanh toán đồng thời (qua biểu đồ hoạt động và biểu đồ tuần tự), đặc tả chi tiết 4 ca sử dụng then chốt cùng các chỉ số yêu cầu phi chức năng.

Trong Chương 3, tôi giới thiệu về nền tảng lý thuyết và công nghệ sử dụng bao gồm cơ sở lý thuyết và lý do lựa chọn bộ công nghệ của đồ án như kiến trúc MERN Stack (MongoDB, Express, React, Node.js), bộ đệm Redis cho rate-limiting và quản lý phiên xác thực, truyền thông thời gian thực Socket.IO, các cổng thanh toán Stripe/VNPay và thuật toán khuyến nghị lai (Hybrid) có suy hao theo thời gian.

Chương 4 tập trung vào phân tích thiết kế, triển khai và đánh giá hệ thống bao gồm thiết kế kiến trúc phần mềm, biểu đồ gói UML, thiết kế lớp chi tiết, thiết kế cơ sở dữ liệu MongoDB, hiện thực hóa các giao diện mua sắm và Dashboard thống kê Recharts dành cho Merchant, cùng với thiết kế các trường hợp kiểm thử chức năng và đánh giá hiệu năng hệ thống.

Trong Chương 5, tôi trình bày các giải pháp và đóng góp nổi bật của đề tài bao gồm cơ chế khóa nguyên tử ngắn hạn dưới cơ sở dữ liệu để kiểm soát lỗi bán vượt mức kho (overselling), lớp bảo vệ phiên xác thực bằng Redis và động cơ tính giá, khuyến mãi xử lý tập trung tại backend.

Chương 6 đưa ra kết luận quan trọng, đánh giá tổng quát mức độ hoàn thành nhiệm vụ và đề tài, những hạn chế còn tồn tại cùng định hướng phát triển, mở rộng hệ thống trong tương lai.

\section*{Kết chương}
Chương 1 đã phác thảo toàn bộ bức tranh tổng quan của đồ án từ bối cảnh bùng nổ của mô hình sàn thương mại điện tử đa nhà bán hàng cho đến những vấn đề cốt lõi về tranh chấp tồn kho đồng thời và an toàn bảo mật hệ thống. Từ các mục tiêu cụ thể, đồ án đã đề xuất định hướng giải pháp toàn diện bằng sự kết hợp giữa kiến trúc MERN Stack, bộ đệm Redis và kênh thời gian thực Socket.IO. Những cơ sở lý thuyết và công nghệ chi tiết làm nền tảng cho việc thiết kế, cài đặt hệ thống này sẽ được thảo luận kỹ lưỡng trong chương tiếp theo – Chương 2.

\end{document}
